% !TEX program = xelatex
% !TEX encoding = UTF-8 Unicode
\documentclass[]{friggeri-cv}

\usepackage{graphicx}
\usepackage{xeCJK}
\setCJKmainfont{Heiti TC}
\usepackage{unicode-math}
\setmathfont{latinmodern-math.otf}

% friggeri-cv's default \section colors only the first 3 tokens of the title (e.g. "sum"+"mary");
% that grabs past short CJK titles like "學歷" and crashes. Color the whole title instead.
\renewcommand{\section}[1]{
  \par\vspace{\parskip}
  {%
    \LARGE\headingfont\color{%
      \ifcase\value{colorCounter}%
        blue\or red\or orange\or green\or purple\else%
        headercolor\fi%
    }#1%
  }
  \stepcounter{colorCounter}
  \par\vspace{\parskip}
}

\begin{document}
\header{孟}{威廉}
       {數據科學 · VR開發 · 統計分析}

\begin{aside}
  \section{聯絡資訊}
    高雄市，台灣
    +886 966 159 200
    \href{mailto:wi.cruzm@gmail.com}{wi.cruzm@gmail.com}
  \section{語言能力}
    西班牙語（母語）
    英語（流利）
    中文（中級）
  \section{程式技能}
    R/RStudio、Python、MATLAB、
    Unity、\LaTeX、VS Code、
    網站開發、MS Office
  \section{核心能力}
    VR應用開發
    資料管理
    技術文件撰寫
    客戶關係管理
\end{aside}

\section{個人簡介}
跨領域專業人士，整合\textbf{認知科學}、\textbf{資料工程}與\textbf{軟體開發}。現為國立成功大學博士候選人，以Unity開發視角依存VR心理旋轉實驗研究空間轉換能力，具備\textbf{R/Python}行為資料分析專長。獨立開發多語言演算法交易平台並發表MQL5指標；三語流利（西班牙語／英語／中文），擅長跨文化技術溝通。

\section{學歷}
\begin{entrylist}
  \entry
    {2019–}
    {心智科學研究所　博士候選人}
    {國立成功大學（NCKU），台南}
    {以視角依存VR心理旋轉任務探討互動性對空間轉換能力之影響。}
  \entry
    {2012–2015}
    {實驗心理學　理學碩士}
    {國立政治大學（NCCU），台北}
    {主修實驗心理學與眼動研究；研究儀器校準與實驗設計。}
  \entry
    {2006–2011}
    {心理學　理學學士}
    {哥倫比亞國立大學，波哥大}
    {主修心理計量與統計分析。}
  \entry
    {2018 / 2015}
    {語言能力檢定證書}
    {TOEFL / TOCFL}
    {英語：101/120（高等）；華語：第三級（中級）。}
\end{entrylist}

\section{工作經歷（一）}
\begin{entrylist}
  \entry
    {08/25–04/26}
    {Live On Language 補習班，高雄}
    {英語教師}
    {針對不同年齡層學生提供英語教學；講授語言結構與溝通技巧。}
  \entry
    {02/25–08/25}
    {Sundex，高雄}
    {國際業務代表暨客戶支援}
    {管理國際客戶帳戶，依客戶規格協調訂單處理與配送。}
  \entry
    {09/18–02/24}
    {劍橋語言學校，台南}
    {西班牙語及英語教師}
    {為西語及英語學習者提供跨文化溝通訓練與語言教學。}
  \entry
    {05/20–02/24}
    {AMC Group}
    {線上英語家教}
    {各程度學員線上英語輔導教學。}
  \entry
    {10/21–02/22}
    {美國鷹語言學校，台南}
    {兼職英語教師}
    {英語初級教學與學習評量。}
  \entry
    {04/16–12/18}
    {ICFES，波哥大}
    {資料管理師}
    {開發全國標準化考試題目資料自動化管理演算法。}
  \entry
    {10/15–04/16}
    {美國鷹語言學校，桃園}
    {英語教師}
    {英語基礎教學。}
\end{entrylist}

\newpage

\section{工作經歷（二）}
\begin{entrylist}
  \entry
    {01/15–09/15}
    {趨勢科技研究部門，台北}
    {研究實習生}
    {垃圾郵件識別演算法開發。}
  \entry
    {09/12–07/15}
    {國立政治大學心理系，台北}
    {研究助理}
    {眼動與閱讀實驗室實驗執行與資料收集。}
  \entry
    {11/11–04/12}
    {哥倫比亞國立大學心理系，波哥大}
    {心理計量師}
    {對受刑人資料進行大規模統計分析，生成風險輪廓。}
  \entry
    {09/11–05/12}
    {安東尼奧·納里尼奧大學心理系，波哥大}
    {研究助理}
    {對教師及學生資料進行統計分析，支援機構審計。}
  \entry
    {03–08/11}
    {哥倫比亞國立大學心理系，波哥大}
    {審計員}
    {全國招募甄試文件審查與驗證。}
\end{entrylist}

\section{著作與研討會}
\begin{entrylist}
  \entry
    {撰稿中}
    {學術論文}
    {國立成功大學}
    {\textit{互動性調節空間轉換能力：來自視角依存VR心理旋轉任務的實證。}}
  \entry
    {05/2023}
    {永續發展目標應用心理學國際研討會暨年會（TAPA），台灣}
    {口頭發表}
    {拖延症是否有解？以自我調節學習為基礎之「拖延心理學」課程設計與學習成效評估；負責統計分析與資料處理。}
  \entry
    {05/2022}
    {視覺科學學會年會（VSS）}
    {壁報發表}
    {利用隨機森林演算法分類輕度腦創傷患者與對照組於視覺空間工作記憶任務中的腦電圖資料。}
  \entry
    {07/2019}
    {第15屆亞太視覺研討會（APCV），大阪}
    {壁報發表}
    {中文句子閱讀中詞彙視覺複雜度分布對眼跳目標之影響。}
  \entry
    {08/2013}
    {第5屆亞洲健康心理學大會，大田}
    {壁報發表}
    {認知行為治療對慢性疲勞症候群及其潛在調節因子效果之統合分析。}
  \entry
    {05/2011}
    {第5屆拉丁美洲健康心理學大會（ALAPSA），哈拉帕}
    {口頭發表}
    {Eurotest 心理計量分析：輕度認知障礙篩查測試。}
  \entry
    {2014}
    {Avendaño-Prieto, B.L., Avendaño-Prieto, G., Cruz-Molina, W., \& Cárdenas-Avendaño, A.}
    {\textit{Diversitas: Perspectivas en Psicología} 10(1), 13--27}
    {非專業研究者多變量統計使用參考指南。}
  \entry
    {2014}
    {Avendaño, B.L., Avendaño, X.A., \& Cruz-Molina, W.}
    {\textit{Avances en Medición} 9, 58--72}
    {Eurotest 失智篩查測試之哥倫比亞版本文化適應研究。}
\end{entrylist}

\clearpage
\section{應用程式與網站}
\begin{entrylist}
  \entry
    {2025}
    {Apex Fintech}
    {台灣}
    {金融服務介面開發。\href{https://apexfintech.net/}{apexfintech.net}}
  \entry
    {2026}
    {Scroll Master}
    {速讀應用程式}
    {協助使用者提升閱讀速度與理解力的速讀網頁應用程式。\href{https://scroll-master.wi-cruzm.workers.dev/}{scroll-master.wi-cruzm.workers.dev}}
  \entry
    {2024}
    {Market Shift \& FVG}
    {MQL5 市場}
    {MetaTrader 5 市場結構與公允價值缺口偵測指標。\href{https://www.mql5.com/en/market/product/131053}{mql5.com}}
\end{entrylist}

\end{document}
